\section*{Bab \ref{ch:g10:proba} --- Peluang dan Penyampelan}

\begin{solution}{exo:g10:proba:1}
Masing-masing dari $6$ hasilnya berpeluang $\frac16$.
``Memperoleh $6$'': $\frac16$.
``Bilangan ganjil'' $= \{1, 3, 5\}$: $\frac36 = \frac12$.
``Paling banyak $4$'' $= \{1,2,3,4\}$: $\frac46 = \frac23$.
``Memperoleh $7$'': \omterm{def:g10:proba:events}{kejadian} mustahil, \omterm{def:g10:proba:distribution}{peluangnya} $0$.
\end{solution}

\begin{solution}{exo:g10:proba:2}
Ada $10$ kelereng yang sama-sama mungkin.
$\P(\text{hijau}) = \frac{5}{10} = \frac12$.
$\P(\text{bukan hitam}) = 1 - \frac{2}{10} = \frac{8}{10} = \frac45$.
Hijau dan kuning \omterm{def:g10:proba:operations}{saling lepas}:
$\P(\text{hijau atau kuning}) = \frac{5}{10} + \frac{3}{10} = \frac{8}{10}
= \frac45$.
\end{solution}

\begin{solution}{exo:g10:proba:3}
Misalkan $S$: ``belajar bahasa Spanyol'', $G$: ``belajar bahasa Jerman''. Maka
$\P(S) = \frac{18}{30}$, $\P(G) = \frac{10}{30}$,
$\P(S \cap G) = \frac{4}{30}$. Aturan penjumlahan:
\[
\P(S \cup G) = \frac{18}{30} + \frac{10}{30} - \frac{4}{30}
= \frac{24}{30} = \frac45 .
\]
\omterm{def:g10:proba:operations}{Komplemennya}: $\P(\text{tidak keduanya}) = 1 - \frac45 = \frac15$.
\end{solution}

\begin{solution}{exo:g10:proba:4}
\omterm{met:g10:proba:tree}{Pohonnya} mempunyai empat jalur, yaitu GG, GA, AG, AA, masing-masing berpeluang
$\frac12 \times \frac12 = \frac14$.
Dua gambar: $\frac14$. Tepat satu gambar (GA atau AG): $\frac24 = \frac12$.
Paling sedikit satu gambar: $1 - \P(\text{AA}) = 1 - \frac14 = \frac34$.
\end{solution}

\begin{solution}{exo:g10:proba:5}
Pengambilan pertama: merah berpeluang $\frac{4}{10}$, putih $\frac{6}{10}$.
Pengambilan kedua \emph{tanpa} pengembalian: setelah merah, tersisa $3$ merah
dan $6$ putih di antara $9$; setelah putih, $4$ merah dan $5$ putih di antara $9$.

Dua merah: $\frac{4}{10} \times \frac39 = \frac{12}{90} = \frac{2}{15}$.

Warna berbeda: jalur MP dan PM:
\[
\frac{4}{10} \times \frac69 + \frac{6}{10} \times \frac49
= \frac{24}{90} + \frac{24}{90} = \frac{48}{90} = \frac{8}{15}.
\]
\end{solution}

\begin{solution}{exo:g10:proba:6}
\emph{1.} Setiap dadu secara saling bebas menampilkan setiap sisi dengan
\omterm{def:g10:proba:distribution}{peluang} $\frac16$, sehingga $36$ pasangan terurutnya sama-sama mungkin; jumlah
matanya ditentukan oleh pasangan itu.

\emph{2.} Jumlah $7$: pasangan $(1,6), (2,5), (3,4), (4,3), (5,2), (6,1)$,
jadi $\P = \frac{6}{36} = \frac16$. Jumlah $12$: hanya $(6,6)$,
$\P = \frac{1}{36}$. Jumlah paling banyak $4$: pasangan yang berjumlah $2$,
$3$ atau $4$: $(1,1)$; $(1,2), (2,1)$; $(1,3), (2,2), (3,1)$ --- enam pasangan,
jadi $\P = \frac{6}{36} = \frac16$.

\emph{3.} Banyaknya pasangan yang mencapai setiap jumlah naik dari $1$ (jumlah
$2$) sampai $6$ (jumlah $7$) lalu menurun: jumlah yang paling mungkin adalah $7$.
\end{solution}

\begin{solution}{exo:g10:proba:7}
Setiap soal dijawab benar dengan \omterm{def:g10:proba:distribution}{peluang} $\frac14$, secara saling
bebas. \omterm{met:g10:proba:tree}{Pohonnya} bertingkat dua:

Keduanya benar: $\frac14 \times \frac14 = \frac{1}{16}$.

Tepat satu benar: jalur benar--salah dan salah--benar:
$\frac14 \times \frac34 + \frac34 \times \frac14 = \frac{6}{16}
= \frac38$.

Tidak satu pun benar: $\frac34 \times \frac34 = \frac{9}{16}$.

Periksa: $\frac{1}{16} + \frac{6}{16} + \frac{9}{16} = 1$.
\end{solution}

\begin{solution}{exo:g10:proba:8}
\emph{1.} \omterm{def:g10:stats:series}{Frekuensi relatif} teramati: $\frac{396}{900} = 0.44$. Dengan
$n = 900$: $\frac{1}{\sqrt{900}} = \frac{1}{30} \approx 0.033$, sehingga
\omterm{def:g10:numbers:interval}{interval} fluktuasinya sekitar $\intcc{0.467}{0.533}$.

\emph{2.} Nilai teramati $0.44$ terletak di luar \omterm{def:g10:numbers:interval}{interval} itu: \omterm{def:g10:proba:sample}{sampel} berisi
$900$ pemilih dari populasi dengan dukungan $50\%$ jarang menyimpang sejauh
itu. Jajak pendapat itu sangat meragukan pengakuan $50\%$ tadi.
\end{solution}

\begin{solution}{exo:g10:proba:9}
\emph{1.} $\P(\text{menang } 6) = \frac16$ (muncul $6$),
$\P(\text{menang } 3) = \frac16$ (muncul $5$),
$\P(\text{menang } 0) = \frac46 = \frac23$.

\emph{2.} Selama $600$ permainan, harapkan sekitar $100$ kali enam, $100$ kali
lima, dan $400$ lemparan lainnya. Kemenangannya: sekitar
$100 \times 6 + 100 \times 3 = 900$. Biayanya: $600 \times 2 = 1200$. Hasil
bersih yang diharapkan: $900 - 1200 = -300$, yaitu rugi sekitar $0.5$ per
permainan secara rata-rata --- permainan itu tidak layak dimainkan.
\end{solution}

\begin{solution}{pb:g10:proba:1}
\textbf{1.} $P(\text{raja atau hati}) = \frac{4}{52} +
\frac{13}{52} - \frac{1}{52} = \frac{16}{52} = \frac{4}{13}$:
raja hati duduk di dalam kedua \omterm{def:g10:proba:events}{kejadian} itu dan tidak boleh terhitung
dua kali --- karena itulah pengurangannya
(\cref{thm:g10:proba:addition}).

\textbf{2.} $P(\text{tanpa gambar}) = \left(\frac12\right)^3 =
\frac18$, sehingga $P(\text{paling sedikit satu}) = \frac78$.

\textbf{3.} $P(MM) = \frac35 \times \frac24 = \frac{3}{10}$;
$P(\text{satu dari tiap warna}) = \frac35 \times \frac24 + \frac25
\times \frac34 = \frac{6}{20} + \frac{6}{20} = \frac35$.

\textbf{4.} $P(BB) = \frac25 \times \frac14 = \frac1{10}$; lalu
$\frac{3}{10} + \frac{6}{10} + \frac{1}{10} = 1$: daun sebuah
\omterm{met:g10:proba:tree}{pohon} selalu berjumlah $1$.

\textbf{5.} \omterm{def:g10:proba:events}{Kejadian} yang \omterm{def:g10:proba:operations}{saling lepas} menuruti $P(A \cup B) = P(A) +
P(B) = 1.1 > 1$: mustahil --- ada \omterm{def:g10:proba:distribution}{peluang} yang dikarang.

\textbf{6.} (Kebanyakan orang, termasuk seribu penulis surat itu,
menjawab ``tidak berpengaruh, 50--50''. Simpanlah jawabanmu untuk
pertanyaan 8.)

\textbf{7.} Mobil di balik pintu 1 (\omterm{def:g10:proba:distribution}{peluangnya} $\frac13$): pemandu acara
membuka salah satu pintu berkambing; berpindah meninggalkan mobilnya:
\emph{kalah}. Mobil di balik pintu 2 ($\frac13$): pemandu acara \emph{harus}
membuka pintu 3; berpindah mendarat pada pintu 2: \emph{menang}. Mobil di
balik pintu 3: setangkup, jadi \emph{menang}. Seluruhnya:
$P(\text{berpindah menang}) = \frac13 + \frac13 = \frac23$.

\textbf{8.} Pilihan pertamamu keliru dengan \omterm{def:g10:proba:distribution}{peluang}
$\frac23$ --- dan tepat pada saat itu, pembukaan yang terpaksa oleh pemandu
acara meninggalkan mobilnya di balik pintu yang tersisa: berpindah mengubah
setiap kekeliruan awal menjadi kemenangan. Jadi $\frac23$, apa pun kata
naluri tadi.

\textbf{9.} Enam skenario yang sama-sama mungkin: mobil di balik $1$, $2$
atau $3$, dan pemandu acara membuka pintu $2$ atau $3$ (tidak pernah
pintumu). Pembukaan itu memperlihatkan mobilnya pada dua di antaranya
(mobil di 2/buka 2, mobil di 3/buka 3): keduanya dibuang. Di antara empat
yang tersisa --- mobil di 1/buka 2, mobil di 1/buka 3, mobil di 2/buka 3,
mobil di 3/buka 2 --- berpindah menang pada dua yang terakhir:
$\frac24 = \frac12$. Dengan pemandu acara yang tidak tahu, 50--50 \emph{memang}
benar: paradoksnya hidup sepenuhnya pada pengetahuan pemandu acara, yang
menyaring skenarionya secara tidak merata.

\textbf{10.} Pilihanmu memuat mobilnya dengan \omterm{def:g10:proba:distribution}{peluang}
$\frac{1}{100}$; $98$ pembukaan oleh pemandu acara yang tahu mengalirkan
sisa $\frac{99}{100}$ itu ke satu pintu tertutup. Berpindahlah,
dan menanglah $99$ kali dari $100$.

\textbf{11.} Tiga kartu (satu bertanda mobil dan dua bertanda
kambing), seorang teman yang mengintip lalu selalu membuang kartu
kambing dari dua kartu yang tidak kamu
pilih; mainkan $300$ ronde, selalu berpindah, lalu catat kemenangannya.
\omterm{def:g10:stats:series}{Frekuensi relatif} yang diharapkan $\frac23 \approx 0.667$, dengan fluktuasi
$\pm\frac{1}{\sqrt{300}} \approx 0.058$: harapkan antara sekitar
$61\,\%$ dan $72\,\%$ kemenangan berpindah --- jauh dari
``50--50'' mana pun.

\textbf{12.} Susunan yang sama-sama mungkin (lebih tua--lebih muda):
LL, LP, PL, PP. ``Paling sedikit satu laki-laki'' menyisakan LL, LP dan PL:
$P(\text{LL}) = \frac13$.

\textbf{13.} ``Yang lebih tua laki-laki'' menyisakan LL, LP:
$P(\text{LL}) = \frac12$. Keterangan yang berbeda menyaring
himpunan skenario yang berbeda --- persis seperti kedua Monty tadi: apa
yang kamu ketahui itu penting, dan \emph{bagaimana} kamu mengetahuinya
sama pentingnya.

\textbf{14.} Ketiganya sama: GGG, AAA: $\frac28 = \frac14$.
Kekeliruan d'Alembert: kedua ``kasusnya'' tidak sama-sama mungkin
--- ``tidak ketiganya sama'' membungkus enam dari delapan hasilnya.
\Cref{prop:g10:proba:uniform} hanya membilang atas hasil yang sama-sama
mungkin.

\textbf{15.} \omterm{def:g10:stats:mean}{Rata-rata} kemenangan per permainan: $6 \times \frac16 + 3
\times \frac16 = 1.50$ euro, berhadapan dengan taruhan $2$ euro: yaitu
rata-rata rugi $0.50$ per permainan, sekitar $300$ euro selama $600$
permainan. Premi yang adil menyamakan rata-rata pembayarannya:
$1\,\% \times 10\,000 = 100$ euro; premi yang sesungguhnya menambahkan biaya
dan margin keamanan --- penjamin asuransi harus bertahan pada tahun yang
buruk, bukan hanya pada tahun yang rata-rata.

\textbf{16.} \omterm{def:g10:numbers:interval}{Intervalnya}: $0.5 \pm \frac{1}{\sqrt{100}} =
\intcc{0.4}{0.6}$. Sebanyak $55$ gambar ($0.55$): di dalam, biasa saja.
Sebanyak $70$ gambar ($0.70$): jauh di luar --- ragukanlah koinnya.

\textbf{17.} $\frac{1}{\sqrt{1000}} \approx 0.032$: seorang calon
yang sebenarnya $50\,\%$ biasanya akan terjaring pada
$\intcc{0.468}{0.532}$ --- dan $0.52$ duduk nyaman di dalamnya.
Jajak pendapat itu \emph{mengisyaratkan} keunggulan tetapi tidak
membuktikan apa pun: batas kesalahannya menelannya.

\textbf{18.} \omterm{def:g10:proba:sample}{Sampel} pertama: \omterm{def:g10:numbers:interval}{intervalnya} $0.04 \pm \frac{1}{20} =
\intcc{0}{0.09}$; nilai teramati $0.07$ ada di dalamnya: tanpa putusan.
\omterm{def:g10:proba:sample}{Sampel} kedua: $0.04 \pm \frac{1}{50} = \intcc{0.02}{0.06}$; nilai
teramati $0.07$ jatuh di luarnya: pengakuannya kini sangat
diragukan. \omterm{def:g10:stats:series}{Frekuensi relatifnya} sama, alatnya lebih tajam: ketelitiannya
tumbuh seperti $\sqrt n$.

\textbf{19.} Fluktuasinya mengecil seperti $\frac{1}{\sqrt n}$:
lipatempatkan \omterm{def:g10:proba:sample}{sampelnya}, deraunya menjadi separuh. Bias --- bertanya
kepada kerumunan yang keliru --- bertahan pada $n$ berapa pun: dua juta
surat suara yang salah pilih tetap saja salah (soal akhir pekan tentang
cara data berdusta di jilid sebelumnya).

\textbf{20.} (i) Periksalah bahwa hasilnya sama-sama mungkin sebelum
dihitung --- d'Alembert lupa. (ii) \omterm{met:g10:proba:tree}{Pohon} mengalikan sepanjang
cabangnya; ``paling sedikit satu'' menyerah kepada \omterm{def:g10:proba:operations}{komplemennya}.
(iii) Keterangan menyaring skenario, dan \emph{sumbernya}
mengubah saringan itu --- Monty yang tahu lawan Monty yang jatuh,
anak sulung laki-laki lawan ada anak laki-laki. (iv) \omterm{def:g10:stats:series}{Frekuensi relatif}
mendekati \omterm{def:g10:proba:distribution}{peluang} dengan laju $\frac{1}{\sqrt n}$, sehingga simulasi
yang sabar menuntaskan perdebatan apa pun --- termasuk perdebatan dengan
sepuluh ribu penulis surat yang marah.

\end{solution}
